\section{A Pause at the Hexad}

There is a natural pause over the Hexad: it represents, in many ways, a repetition of the Monad, being halfway in between 1 \& 10. The Hexad, or Seal of Solomon, reminds us that even the Monad itself is not original, but ab-original, because Iamblichus never deals with the Number Zero. We might wonder if Zero represents Non-Being, something which Greek civilization did not investigate as thoroughly as it did the Circle and its Center as The Beginning. A bright mathematician and philosopher/poet might add to Iamblichus' meditations by exploring what the East discerned Non-Being to be, and how this might be integrated into Western metaphysics through Iamblichus' meditative sequence on Number.

Thus, when the Hexad begins the surge of re-integrative Life, we are reminded that the One or Logos as we know it in the West is itself not God, in a sense, because the chain of progression makes us consider the subtle variations between numbers (as well as Zero). As Jesus said, Why do you call me Good? There is One Good, who is God.'' It is no coincidence that Islam arose between East and West, in the wastes of the desert, where men had not forgotten that the scandal of the Incarnation was a fire which might burn those who neglected the study of Non-Being. In fact, the Orthodox argue that Catholic mysticism, over time, degenerated into emotionalism and sentimentality of a sort, which is related to the Western emphasis on the second Chakra (the Logos). As this author points out\footnote{\url{http://paulijungunusmundus.eu/rfr/radbilde.htm}}, this drive lead them to abandon the Supreme Eros (the front Chakra) and also meditation in favor of a degenerated concept of Logos, which leads to Bacon and modern empiricism\footnote{\url{http://www.gornahoor.net/?p=3468}}. The alchemical process degenerated into the world of the Machine. Yet even here, there are echoes\footnote{\url{https://en.wikipedia.org/wiki/Quark\_model}}, intimations of ``deprivation'' as George Parkin Grant termed them -- we learn the existence of Heaven by experiencing Hell. Even raw physics teaches the importance of the Hexad.

In a future essay, we will see how Tolkien himself (the last man to take up arms against the Machine in letters, as a crusade) anticipated this in his myth of the Lost Road\footnote{\url{https://en.wikipedia.org/wiki/The\_Lost\_Road\_and\_Other\_Writings}}. The world no longer had a straight road ``into the West'' that rose into the upper air, but became Round. The road to Non-Manifestation and Non-Being was closed to the Western Man.

If we recover a doctrine of the Hexad, we might balance not only the Monad and the end of the cycle in 10, but also see how Zero teaches us even more about Origins. (Please note: in a sense, Iamblichus' Monad can be said to contain both the 0 and the 1, and so is not incomplete -- rather, we say that Iamblichus was more concerned with Manifestation than Non-Being).

The Hexad is the Seal of Solomon, which the Templars were created to embody. Here was yet another opportunity afforded by the far reach of Heaven to remind the West of what it should already know. Some mystics like Nicholas of Flue realized that there had to be two Trinities, since the lower realms had to be constantly related back to, not only the metaphysical principles of realization in the manifested Trinity, but united in what alchemy was designed to reach -- perfect reintegration into the Monad of the Monad. This was to occur in the Heart, where was hid all the secrets of earth and heaven. Even in Islam, which kept the echoes alive, a need was seen by the Sufis to recognize the truth which Christianity held in potential, and the heart remained paramount to exoteric doctrine focused on the One.

The mystery of the Hexad, or double Trinity, teaches man that God is a God of the living, and that even correct dogma or doctrine was not everything -- there remained the task of the Heart, which Valentin Tomberg undertook to explicate in his Meditations.

The Hexad re-unites Wood (the mortal Fifth Element) with the true Fifth Element or Quintessence: the ``Hearth'' (another name for the Monad) is in between Earth and Heaven, just as the Heart mediates the lower 3 chakras with the upper 3. God must swear a double oath -- one by Earth, one by Heaven (see the Epistle to the Hebrews). This is one of the true meanings of the Incarnation, which was vouchsafed to the Christians, if they can remember it.


\flrightit{Posted on 2013-02-15 by Logres}

\begin{center}* * *\end{center}

\begin{footnotesize}
\begin{sffamily}
\texttt{scardanelli on 2013-02-16 at 11:22 said:}

I'm eager to read your analysis of Tolkien's Lost Road myth. There's an interesting article here by Joscelyn Godwin discussing the Traditional nature of Tolkien's work:

\url{http://hermetic.com/godwin/tolkien-and-the-primordial-tradition.html}


I've been thinking lately that perhaps the creation of new storytelling and myths might be the only way to reach much of the Western population; a way of indirectly planting the ``seed'' of traditionalism without invoking the reaction that the modern world has trained into them. Tolkien could be a model for this work, as he has obviously mastered the genre.

\hfill

\texttt{Caleb Cooper on 2013-02-16 at 19:42 said:}

I think just as you can say the Dyad is like the Monad to the levels below it, the Monad is also like the Dyad to the Zero. Perhaps there's a general principal in play here that the proper order of things is two-faced; the feminine reflective principle faces towards the superior level before it, while the masculine aspect faces the other way, radiating to the lower levels what has been assimilated from above. We can see examples of this in Jesus' complete submission to the Father; the Church as the bride of Christ, ect... 

\hfill

\texttt{Logres on 2013-02-17 at 06:34 said:}

Caleb -- you nailed it: you just articulated what has been lurking in my mind.

Scardanelli -- you are also spot on : that was the whole mission of the Inklings, in which they failed (for several reasons), but their instinct was sound, and it might work, in fact stands a good strong chance of doing just that. We'll look at how Tolkien (who came closest in terms of poetic power and natural instinct) started this. I am not sure if Williams had that much esoteric influence (Tolkien certainly had more popular influence): I think Tolkien distrusted his influence over Lewis. Tolkien had more talent, and Barfield had to work, as well as re-invent the wheel, to a degree.

\hfill

\texttt{Rui Artur on 2022-08-11 at 11:36 said:}

The passage from Zero to One (Beyond Being to Being) is one of the things I read in the Apocalypse verse `the lamb was slain from the foundation of the world', because Zero cannot create (every number multiplied by it is Zero), and exactly the opposite is the case with One, which multiplied by any number produces that number (the ten thousand things of the Tao Te Ching). So the Zero had to `go out of himself', had to be `sacrificed' to produce things outside it.


\hfill

\end{sffamily}
\end{footnotesize}

